% ===== ch06 : 原文 main.tex 第 1301--1478 行（4a 后半 / 4b 商空间的缠绕 Fukaya 范畴 / 4c Künneth / 4d Kartal 修改）=====
% 本文件是原文 LaTeX 源码的逐行分片（去掉整段注释掉的旧稿），仅供翻译时对照；最终译文写到 src/chunks/ch06.tex。
        \begin{scope}[yshift=-3cm,xshift=4cm]
            \node at (0.5,1) {$\vdots$};
        \end{scope}
        \end{tikzpicture}
    \caption{$\mathbb{Z}^2$-cover of the Twice-Punctured Disk with Stops.}\label{fig:Z2-cover-twice-punctured-disk}
    \end{figure}
    As we know the generation of each sectors in this gluing, using sectorial descent for the diagram \[
        \begin{tikzcd}
            \mathcal{W} \left(\bigsqcup_{\mathbb{Z}^2}T^{\ast}[0,1]\right)\arrow[r]\arrow[d]&\mathcal{W} \left(\bigsqcup_{\mathbb{Z}}(\mathbb{C},e^{z_1}+e^{-z_1})\right)\\ 
            \mathcal{W}\left(\bigsqcup_{\mathbb{Z}}(\mathbb{C},e^{z_2}+e^{-z_2})\right)&
        \end{tikzcd}
    \]we see that $\mathcal{X}$ is still generated by the union of generators of the Fukaya categories \[\mathcal{W}\left(\bigsqcup_{\mathbb{Z}} (\mathbb{C},e^{z_i}+e^{-z_i})\right),\]which is again invariant under the $\mathbb{Z}^2$-action.\par
    Based on this gluing diagram, we can label generating Lagrangian submanifolds of $\mathcal{W}(\mathcal{X})$ as $L_i [a,b](c)$, where $i\in\{0,1\}$ denotes whether this Larangian is from the left-hand punctured disk or the right, $c\in\{0,1\}$ denotes the ``slope'' of $L_i$ and $(a,b)\in\mathbb{Z}^2$ denotes which connected component and after assigning a canonical $0$ which sheets in the $\mathbb{Z}$-cover it lies in.\par
    To better understand the Floer cohomologies $CW^{\ast}\left(L_i [a,b](c),L_j[c,d](c);\Bbbk\right)$, it's more intuitive to look at the twice-punctured disk model and use imagination for where the chosen Lagrangian gets lifted. See figure \ref{fig:computation-Z2-cover}. Suppose that $i=j$, as the wrapping around the middle circle clockwise shifts the bi-degree by $(-1,1)$, we see that \[
        CW^{\ast}\left(L_i [a_1,b_1](c_1), L_i [a_2,b_2](c_2)\right)\simeq\left\{\begin{array}{ll}
            \Bbbk,&\text{if}\ \vert a_1 -a_2\vert\leq 1,\ b_1= b_2\ \text{and}\ c_2-c_1=1;\\ 
            \Bbbk,&\text{if}\ a_1-b_1=b_2-a_2\ \text{and}\ c_1=c_2;\\ 
            0,&\text{otherwise}.
        \end{array}\right.
% ---- 原文行 1320 ----
    \]
    \begin{figure}[ht]
        \centering
        \begin{tikzpicture}
            \begin{scope}
        \draw [thick] (-2,0) arc[start angle=0, end angle=360, x radius=0.5, y radius=1];
        \draw [thick] (-2.5,1) -- (0.5,1);
        \draw [thick] (-2.5,-1) -- (0.5,-1);
        \draw [thick,red!50!black] (0.5,1) arc[start angle=90, end angle=-90, x radius=0.5, y radius=1];
        \draw [thick,red!50!black, dashed] (0.5,-1) arc[start angle=270, end angle=90, x radius=0.5, y radius=1];
        \draw [thick,blue!50!black] (-2,0) .. controls (-1,0) .. (0.5,1) node [pos=0.5,below] {$L_1[a,b](0)$};
        \draw [thick,blue!50!black] (-2.2,{sqrt(1-(0.3/0.5)^2)}) .. controls (-1.5,{sqrt(1-(0.3/0.5)^2)}) .. (-1,1) node [above] {$L_1[a',b'](1)$};
        \draw [thick,dashed,blue!50!black] (-1,1) .. controls (-0.8,-1) .. (-0.5,-1);
        \draw [thick,blue!50!black] (-0.5,-1) .. controls (0.5,-1) .. (0.5,1);
        \end{scope}
        \node [ellipse,
    minimum width=3mm,
    minimum height=1.5mm,
    inner sep=0pt,
    fill=red!50!black,
% ---- 原文行 1340 ----
    draw=red!50!black,
    thick] at (-2.5,-1) {};
    \begin{scope}[xshift=3cm]
        \draw [thick] (-2.5,1) -- (0.5,1);
        \draw [thick] (-2.5,-1) -- (0.5,-1);
        \draw [thick] (0.5,1) arc[start angle=90, end angle=-90, x radius=0.5, y radius=1];
        \draw [thick, dashed] (0.5,-1) arc[start angle=270, end angle=90, x radius=0.5, y radius=1];
        \node [ellipse,
    minimum width=3mm,
    minimum height=1.5mm,
    inner sep=0pt,
    fill=red!50!black,
    draw=red!50!black,
    thick] at (0.5,-1) {};
    \draw [thick,blue!50!black] (1,0) .. controls (0,0) .. (-2.5,1) node [pos=0.5,below] {$L_2[a,b](0)$};
    \end{scope}
    \node [ellipse,
    minimum width=3mm,
    minimum height=1.5mm,
    inner sep=0pt,
% ---- 原文行 1360 ----
    fill=white,
    draw=red!50!black,
    thick] at (0.5,1) {};
        \end{tikzpicture}
        \caption{Computing wrapped Floer cohomology from the base space.\label{fig:computation-Z2-cover}}
    \end{figure}
    Similarly, when $i\neq j$, it can be read from figures \ref{fig:Z2-cover-twice-punctured-disk}, \ref{fig:computation-Z2-cover} that \[
        CW^{\ast}\left(L_i [a_0, b_0](c_i),L_j[a_1,b_1](c_j);\Bbbk\right)=\left\{\begin{array}{ll}
            \Bbbk,&\text{if}\ a_1-b_1 =b_0-a_0;\\ 
            0,&\text{otherwise}.
        \end{array}\right.
    \]This provides a complete characterization of the wrapped Fukaya category. The quiver presentation will consist of a 3D graph, and is omitted here. Note that some parallel description of the same space also appears in \cite{Neumann04} in the study of extended Bloch groups.
\end{exmp}
We end this section with a discussion of a higher-dimension situation.
Recall that for a smooth manifold $M$, with a choice of a Morse function $f\colon M\to\mathbb{R}$ on $M$ and a Riemannian metric $g$ the cotangent bundle $T^{\ast} M$ carries a natural Weinstein structure $(\phi,\lambda_f,\omega_f)$ where $\phi (x,\alpha) =f(x)+\vert\alpha\vert_g^2$, $\lambda_f =df-\lambda_{\can}$ with $\lambda_{\can}$ the canonical Liouville structure on $T^{\ast} M$ given by $\pi^{\ast}\lambda_{\can} =\lambda_{\can}$ where $\pi\colon T^{\ast} M\to M$ is the bundle map, and $\omega_f =d\lambda_f$.\par
A Legendrian embedding $\Lambda\hookrightarrow\partial_{\infty} T^{\ast} M$ provides a symplectic hypersurface $T^{\ast}\Lambda\hookrightarrow T^{\ast} M$, which itself carries a Weinstein structure as above.
Due to \cite{Eliashberg18}, after modifying the Weinstein structure on $T^{\ast} M$ we can get a Weinstein structure making $T^{\ast}\Lambda\subseteq T^{\ast} M$ a Weinstein hypersurface, defining a Weinstein sector due to \cite{GPS20}.\par
One example appears in \cite{FLTZ12}, where $M=\mathbb{T}^n$ and for a given polygon $\Sigma$ of dimension $n$ we can construct a Legendrian link $\Lambda_{\Sigma}\hookrightarrow\partial_{\infty} T^{\ast}\mathbb{T}^n$ to obtain a Weinstein sector $(T^{\ast}\mathbb{T}^n,\Lambda_{\Sigma})$.
The $\mathbb{Z}^n$-cover $(T^{\ast}\mathbb{R}^n,\widetilde{\Lambda}_{\Sigma})$ with the pull-back Weinstein structure is also naturally a Weinstein sector, by our definition.\par
Note that $\pi_1 (T^{\ast}\mathbb{T}^n,\ast)\cong\mathbb{Z}^n$, hence in the $\mathbb{Z}^n$-cover the $1$-handle attachments are altered, while all remaining handles are lifted trivially.
% ---- 原文行 1380 ----
Write $X=(T^{\ast}\mathbb{R}^n,\widetilde{\Lambda}_{\Sigma})$, we can then construct an exhausting sequence of Weinstein cobordisms $X_0\hookrightarrow X_1\hookrightarrow X_2\hookrightarrow\dotsb\hookrightarrow X$ as follows.\par
$X_0$ is a disjoint union of finitely many points, with labels $0$ or $e_i$ where
$$e_i =(0,\dotsb ,0,\underbrace{1}_i,0,\dotsb ,0).$$
$X_1$ is obtained from $X_0$ by attaching $1$-handles to $X_0$ connecting the point $0$ to all the $e_i$, all the remaining handles of higher index lying over these $1$-handles, and half-handles attached to the given handles that contribute to the Legendrian stops near the contact boundary-at-infinity.
$X_2$ is constructed from $X_1$ by adding $0$-handles with labels $e_i+e_j$, $1\leq i,j\leq n$, and repeating the procedure of attaching handles described in the construction of $X_1$.
This procedure can then be continued and we will exhaust all the handles for the lifted Weinstein manifold $(T^{\ast}\mathbb{R}^n,\widetilde{\Lambda}_{\Sigma})$.
It's immediate from the definition that
\begin{lmm}\label{lmm:Fukaya-category-Gmn-equivariant-generation}
    The wrapped Fukaya category $\mathcal{W} (T^{\ast}\mathbb{R}^n,\widehat{\Lambda}_{\Sigma})$ is generated by cotangent fibres $T^{\ast}_q\mathbb{R}^n$ where $q\in\mathbb{Z}^n$ and all linking disks to the (infinite type) Legendrian link $\Lambda_{\Sigma}$ near infinity.
\end{lmm}
On the other hand, we can construct from the data $(T^{\ast}M,\Lambda)$ dg-category of constructible sheaves of categories $\Sh_{\Lambda} (M)$ with singular supports contained in the Whitney stratification $\mathcal{S}_{\Lambda}$ of conic isotropic submanifolds consisting of the zero section together with positive real cones over the Legendrians $\Lambda$ near infinity valued in the dg-category of chain complexes over $\Bbbk$ \cite{GPS24M}.
It's proved in \cite{GPS24M} that
\begin{thm}[Sheaf-Floer correspondence]\label{thm:sheaf-Floer-correspondence}
    Suppose that $M$ is a smooth manifold with $f\colon M\to\mathbb{R}$ a Morse function with finitely many critical points. There is a canonical quasi-equivalence of categories $$
        \Perf\mathcal{W} (T^{\ast} M,\Lambda)^{op}\simeq\Sh_{\Lambda} (M)^{\omega}
    $$where linking disk at a point $p\in\Lambda$ are sent to co-representatives of microstalk functor associated to $p$ and cotangent fibres $T^{\ast}_p M$ with $p$ not in the image of the projection of $\Lambda$ under $\partial_{\infty} T^{\ast} M\to M$ sent to a co-representative of the stalk functor at $p$.
\end{thm}
We will provide a generalization of this theorem to infinite-type sectors in this specific situation.
\begin{prp}\label{prp:inf-type-sheaf-Floer-Rn}
    For $M=\mathbb{R}^n$ and $\widetilde{\Lambda}_{\Sigma}$ the lifted FLTZ skeleton, there is a quasi-equivalence of $A_{\infty}$-categories \[
% ---- 原文行 1400 ----
        \Perf\mathcal{W} (T^{\ast}\mathbb{R}^n,\widetilde{\Lambda}_{\Sigma})^{op}\simeq\Sh_{\widetilde{\Lambda}_{\Sigma}} (\mathbb{R}^n)^{\omega}.
    \]
\end{prp}
\begin{proof}
    Note that the exhaustion $X_0\hookrightarrow X_1\hookrightarrow\dotsb\hookrightarrow X$ by construction are exhaustions by cotangent bundles with Legendrian stops, and therefore for each $X_i=(T^{\ast} M_i,\Lambda_i)$, we have by \cite{GPS24M} a canonical quasi-equivalence of categories \[
        \Perf\mathcal{W} (X_i)^{op}\simeq\Sh_{\Lambda_i} (M_i)^{\omega}.
    \]Under this quasi-equivalence, Viterbo functor is given by the restriction functor \[
    \Sh_{\Lambda_{i+1}} (M_{i+1})\to\Sh_{\Lambda_i} (M_i)
    \]sending a constrctible sheaf to its restriction to an open subset removing some of the strata, which sends compact objects to compact objects using the description \cite[Lemma 4.18]{GPS24M}.
    We can then verify directly using the exhaustion that we can present $\Sh_{\widetilde{\Lambda}_{\Sigma}} (\mathbb{R}^n)$ as a homotopy limit \[
        \Sh_{\widetilde{\Lambda}_{\Sigma}} (\mathbb{R}^n)\simeq\holim\left(\Sh_{\Lambda_0} (M_0)\leftarrow\Sh_{\Lambda_1} (M_1)\leftarrow\dotsb\right).
    \]\cite[Proposition 4.20]{GPS24M} together with the definition of wrapped Fukaya categories then implies the desired quasi-equivalence.
\end{proof}
\subsection{Wrapped Fukaya category of quotients.} Let $L,K\subseteq X/G$ be exact cylindrical spin Lagrangian submanifolds such that $\langle\pi_1 (L),\pi_1 (K)\rangle\cap G= \{1\}$ in $\pi_1 (X/G)$, then we have
\begin{prp}
    $CF^{\ast} (L,K) =\displaystyle\bigoplus_{g\in G} CF^{\ast} (g(\tilde{L}),\tilde{K})$ as graded $\Bbbk$-modules. Moreover, for given positive isotopy $\psi\colon K\leadsto K^+$ with lift $\tilde{\psi}$, the continuation map $c_{\psi}\colon CF^{\ast} (L,K)\to CF^{\ast} (L,K^+)$ coincides with $\displaystyle\bigoplus_{g\in G} c_{\tilde{\psi}}\colon \bigoplus_{g\in G} CF^{\ast} (g(\tilde{L}),\tilde{K})\to\bigoplus_{g\in G} CF^{\ast} (g(\tilde{L}),\tilde{K}^+)$.
\end{prp}
\begin{proof}
    An intersection point $p\in L\cap K$ in $X/G$ corresponds to a point $\tilde{L}\cap g(\tilde{K})$ in $X$ for some $g\in G$, and therefore by definition of Floer cochain complex we have $CF^{\ast} (L,K)=\displaystyle\bigoplus_{g\in G} (g(\tilde{L}), \tilde{K})$. Given a positive isotopy $\psi\colon K\leadsto K^+$ with a given lift $\tilde{K}$, using the fact that $X\to X/G$ is a covering space we can lift the whole isotopy $\psi$ to an isotopy $\tilde{\psi}\colon \tilde{K}\leadsto\tilde{K}^+$. Under the $G$-action we have disjoint copies of isotopies $g(\tilde{\psi})\colon g(\tilde{K})\leadsto g(\tilde{K}^+)$, and therefore provides Floer-theoretic continuation maps \[
        c_{g(\tilde{\psi})}\colon CF^{\ast} (g(\tilde{L}),\tilde{K})\to CF^{\ast} (g(\tilde{L}),\tilde{K}^+)
% ---- 原文行 1420 ----
    \]for each chosen Lagrangian $\tilde{L}$ in $X$. Taking the direct sum provides the desired continuation map.
\end{proof}
Now given a sequence of admissible Lagrangian submanifolds $L_0, L_1,\dotsb ,L_k$ in $X/G$ with intersection points $p_i\in CF^{\ast} (L_{i-1}, L_i)$ for $1\leq i\leq k$ and $p_0\in L_k\cap L_0$, a regular $J$-holomorphic disk $u\in\mathcal{M} (p_0;p_1,\dotsb ,p_k)$ admits a unique natural lift to an element \[
\tilde{u}\in\mathcal{M}(X;\tilde{p}_0;\tilde{p}_1,\dotsb ,\tilde{p}_k)
\]where for fixed chosen lift $\tilde{L}_i$ of $L_i$, the intersection points $\tilde{p}_i\in g_i(\tilde{L}_{i-1})\cap\tilde{L}_i$ lift to intersection points of $g_i (\tilde{L}_{i-1})$ and $\tilde{L}_i$ for some $g_i$ modulo $k$, and hence the lifted holomorphic disk $\tilde{u}$ contributes to a map \[
\mu^k\colon CF^{\ast} (g_{k-1}(\tilde{L}_{k-1}),\tilde{L}_k)\otimes\dotsb\otimes CF^{\ast} (g_1\dotsb g_{k-1}(\tilde{L}_0),g_2\dotsb g_{k-1} (\tilde{L}_1))\to CF^{\ast} (g_1\dotsb g_k (\tilde{L}_0),\tilde{L}_k)[2-k].
\]
Conversely, any regular $J$-holomorphic disk in $u$ with boundary on $\displaystyle\bigsqcup_{g\in G}g(\tilde{L})$ asymptotic to intersections $g(\tilde{L}_{i-1})\cap\tilde{L}_i$ projects to regular $J$-holomorphic disks in $X/G$ with boundaries on $L_i$, and therefore the $A_{\infty}$-structure on $\mathcal{W} (X/G)$ can be described as $A_{\infty}$-structures of $\mathcal{W} (X)$ that is invariant under the $G$-action.\footnote{The same argument was also used in Seidel's work \cite{Seidel11}.} This is enough to conclude the theorem.
\begin{proof}[Proof of Theorem \ref{thm:wrapped-Fukaya-quotient}]
    Note from definition \ref{defn:wrapped-Fukaya-cat-inf-sectors} that the wrapped Fukaya category $\mathcal{W} (X)$ is generated by all the cocores appearing in any of the exhausting sequence of Weinstein domains. In particular, the $G$-invariant Weinstein structure on $X$ provides $G$-copies of critical points associating to a single critical point of any given Weinstein structure on $X/G$. Therefore we conclude that the disjoint union $\displaystyle\bigsqcup_{g\in G} g(\tilde{L})$ of cocores, where $\tilde{L}$ is a lift of cocore of the Weinstein structure on $X/G$, generates the wrapped Fukaya category of $X$. On the chain level we have \[
        CW^{\ast}_{X/G}\left(L,K;\Bbbk\right)=\bigoplus_{g\in G} CW^{\ast}_X (g(\tilde{L}),\tilde{K};\Bbbk),
    \]where the $A_{\infty}$-structures are written as the discussion above, which pairs exactly with the definition of $A_{\infty}$-orbit categories in Definition \ref{defn:Ainf-orbit-category}.
\end{proof}
\subsection{The K\"unneth theorem and deformation invariance.} Recall from \cite{GPS24S} that we have the following K\"unneth theorem for wrapped Fukaya categories:
\begin{thm}[{\cite[Theorem 1.5]{GPS24S}}]\label{thm:Kunneth-theorem-Fukaya-cat}
    For Weinstein sectors $X$ and $Y$ of finite type, there is a quasi-equivalence of $A_{\infty}$-categories \[
        \mathcal{W} (X)\otimes\mathcal{W} (Y)\simeq\mathcal{W} (X\times Y).
    \]
\end{thm}
The result we need to use in our situation is slightly more general: it allows $X$ to be of infinite type, and so $\mathcal{W} (X)$ is defined instead in an ad hoc way. Nonetheless, we can still achieve the following:
% ---- 原文行 1440 ----
\begin{crl}\label{crl:Kunneth-Fukaya-infinite}
    Suppose that $X$ is a Weinstein sector and $Y$ is another Weinstein sector of finite type, then there is a quasi-equivalence of $A_{\infty}$-categories \[
        \mathcal{W} (X)\otimes\mathcal{W} (Y)\simeq\mathcal{W} (X\times Y).
    \]
\end{crl}
\begin{proof}
    Let $\emptyset=X_{-1}\subseteq X_0\subseteq\dotsb\subseteq X$ be an exhaustion of $X$ by Weinstein domains, then for each $n$ we have quasi-equivalences of $A_{\infty}$-categories \[
        \mathcal{W} (X_n)\otimes\mathcal{W} (Y)\xrightarrow{\simeq}\mathcal{W} (X_n\times Y).
    \]Note that Ornaghi's tensor product \cite{Ornaghi25} is functorial, and so the Viterbo restriction $\mathcal{W} (X_{n+1})\twoheadrightarrow\mathcal{W} (X_n)$ gives rise to $A_{\infty}$-functors \[
        \mathcal{W} (X_{n+1})\otimes\mathcal{W} (Y)\twoheadrightarrow\mathcal{W} (X_n)\otimes\mathcal{W} (Y)
    \]which fits into a homotopically commutative diagram of $A_{\infty}$-categories \[
        \begin{tikzcd}
            \mathcal{W} (X_{n+1})\otimes \mathcal{W} (Y)\arrow[r,"\simeq"]&\mathcal{W} (X_{n+1}\times Y)\\ 
            \mathcal{W} (X_n)\otimes\mathcal{W} (Y)\arrow[u,two heads]\arrow[r,"\simeq"]&\mathcal{W} (X_n\times Y)\arrow[u,two heads]
        \end{tikzcd}
    \]Note that $\{X_n\times Y\}_n$ is an exhaustion of $X\times Y$ by Weinstein subdomains, and therefore by taking homotopy limits we get the quasi-equivalence \[
        \mathcal{W}^{big} (X)\otimes\mathcal{W}(Y)\simeq\mathcal{W}^{big} (X\times Y),
    \]and by keeping track of the ad-hoc definition of \ref{defn:wrapped-Fukaya-cat-inf-sectors} we get the desired K\"unneth isomorphism.
\end{proof}
The other essential ingredient is a proof that certain exact symplectic deformation of an infinite-type Weinstein sector does not change the corresponding Fukaya category.
% ---- 原文行 1460 ----
\begin{prp}\label{prp:deformation-invariance-Fukaya}
    Suppose that $X$ is a Weinstein sector with Weinstein structure $\phi$ and $\{\phi_t\}$ a Weinstein homotopy with $\phi_0=\phi$ such that for each $n$, zeroes of $\phi_t\vert_{X_n}$ are contained in a compact subset in the interior of $X_n$. Then $\mathcal{W} (X,\phi_0)\simeq\mathcal{W} (X,\phi_1)$.
\end{prp}
\begin{proof}
    From Lemma 2 of \cite{Jeffs22} we know that $\mathcal{W} (X_n,\phi)\simeq\mathcal{W} (X_n,\phi_1)$ for each $n$, and by definition of the Weinstein homotopy we see that this quasi-isomorphism is compatible with the Viterbo restriction. Therefore by taking the homotopy limit we get the desired deformation invariance.
\end{proof}
\subsection{Kartal's modification.} We end this section with a recall of Kartal's modification \cite{Kartal21} adapted to the current situation. Let $X$ be a Weinstein sector of finite type with Weinstein structure $(\varphi,\lambda)$ and an exact symplectomorphism $\phi\colon X\to X$ that is contact type in a cylindrical end of $X$, we are going to construct a Weinstein structure $\psi$ on $(T^{\ast}\mathbb{R}\times X,e^z+e^{-z})$ such that $\psi$ is invariant under the $\tau^{-1}\times\phi$-action, with $\tau\colon (T^{\ast}\mathbb{R},e^z+e^{-z})\to (T^{\ast}\mathbb{R},e^z+e^{-z})$ the translation-by-$1$ symplectomorphism.\par
The ansatz, from \cite{Kartal21}, is to consider Weinstein structures of the form\[
\pi_X^{\ast}\varphi_z +p^{\ast}\varphi_{(T^{\ast}\mathbb{R},e^z+e^{-z})}
\]where $\{\varphi_z\}$ is a $T^{\ast}\mathbb{R}$-family of Weinstein structures on $X$ so that $\phi^{\ast}\varphi_{z+1}=\varphi_z$ and $p\colon (T^{\ast}\mathbb{R}\times X,e^z+e^{-z})\to (T^{\ast}\mathbb{R},e^z+e^{-z})$ is the projection to the $(T^{\ast}\mathbb{R},e^z+e^{-z})$-factor. We can further assume that $\varphi_{z=x+iy}$ is $y$-invariant and is constant when $x\in\mathbb{Z}$ and so the resulting Weinstein deformation does not change the stop and we can simply write $\{\varphi_s\}_{s\in\mathbb{R}}$. By definition of $\phi$ we are able to find a smooth function $K\colon X\to\mathbb{R}$ where $\phi^{\ast}\lambda_X=\lambda_X+dK$, and the construction of $\{\lambda_s\}$ is given by considering a smooth family of functions $\{K_s\}_{s\in\mathbb{R}}$ that is cylindrical in a cylindrical end of $X$ so that $\lambda_s =\lambda +dK_s$. The condition $\phi^{\ast}\lambda_{s+1}=\lambda_s$ then turns into the condition that \[
    dK_s=d(\phi^{\ast} K_{s+1}+K).
\]We can then find a smooth family of functions $K_s$ such that $\phi^{\ast} K_{s+1} +K-K_s=0$ holds for all $s\in\mathbb{R}$, by exactly the same procedure as described in the proof of Proposition 2.1 in \cite{Kartal21}.\par
To upgrade this Liouville structure to a Weinstein structure, we need to further choose a family of smooth functions $\varphi_s$ together with a constant $C>0$ so that $\varphi_C =p^{\ast}\varphi_{(T^{\ast}\mathbb{R},e^z+e^{-z})} +C^{-1}\varphi_z$ is Lyapunov for the Liouville structure $\lambda_C =p^{\ast}\lambda_{(T^{\ast}\mathbb{R},e^z+e^{-z})} +C^{-1}\lambda_z$. The conclusions drawn in \cite{Kartal21} is then 
\begin{prp}[{Similar to Proposition 2.3 of \cite{Kartal21}}]\label{prp:Kartal-modification-Weinstein}
    There exists a Weinstein structure on $(T^{\ast}\mathbb{R}\times X,e^z+e^{-z})$ of the form \[
        \varphi_{(T^{\ast}\mathbb{R}\times X,e^z+e^{-z})} =p^{\ast}\varphi_{(T^{\ast}\mathbb{R},e^z+e^{-z})}+C^{-1}\varphi_z,\quad \lambda_{(T^{\ast}\mathbb{R}\times X,e^z+e^{-z})}=\pi^{\ast}\lambda_{(T^{\ast}\mathbb{R},e^z+e^{-z})} +C^{-1}\lambda_z
    \]for $C>0$ sufficiently large, that is invariant under the $\tau^{-1}\times\phi$-action on $(T^{\ast}\mathbb{R}\times X,e^z+e^{-z})$.
\end{prp}
% \newpage
